\subsection{The E6 selected cusp with its actual finite cap}
\begin{statement}{P3K.95}{Proved}
\par\medskip\noindent\textbf{Statement.}\par\nopagebreak[4]

For the actual source family
\[
 Y^2=X^3+(3-2w^2)X^2+(w^2-3)TX+T^2,\qquad
 G=(0,T),\qquad q=T-(w^2-1),
\]
retain \(M_6=\mathcal O(6G-O)\), source suspension by \(-G\),
the complete selected-zero Hom section, and the quadratic reduced-IV
cap with either primitive norm sign \(n=\pm1\), its prescribed
finite character and full integral annular reference cochain.
Fix the sufficiently small nonzero scalar, including its chosen
root and phase. The actual positive \(w\)-loop about zero has,
on every selected lattice-basis fan, the same oriented sphere return
as the complete K71 selected-cusp algorithm. This equality is
induced by the full source, Cartier, finite-root, exterior, normal
and Cox maps below.

On the standard vertical fan \(T_j\), \(j\in\mathbb Z\), its
physical characteristic pair and oriented class are
\[
 (k_M,\ell_M)=(-6-18j(j-1),0),\qquad
 [\Sigma_j]=-1-3j(j-1)\pmod {28}.
\]
The corresponding native free pair is
\((0,-6-18j(j-1))\). At \(j=0\) the complete native triple is
\((0,-6,-1)\). For every other selected lattice-basis fan the
actual unit matrices and native trace maps in EC.6 give the same
K71 calculation. The marked parameter comparison to the original
source is
\[
 r(w)=\frac{3-2w^2}{(w^2-3)^2},\qquad
 r(w)-\frac13=-\frac{w^4}{3(w^2-3)^2};
\]
the \(w\)-loop therefore has actual local degree four over the
\(r-\tfrac13\) parameter. All conclusions refer to this specified
ordered loop and the retained \(M_6,\ b=1,\ m=2\) completion.

\end{statement}
\begin{proof}
\par\medskip\noindent\textbf{EC.1. The full physical-height equation and the ordered normalization.}\par\nopagebreak[4]

Set
\[
 \begin{gathered}
 v=w^2-1,\quad D=w^2-3,\quad G=(0,T),\quad q=T-v,\\
 E_w:\ Y^2=X^3+(3-2w^2)X^2+DTX+T^2 .
 \end{gathered}
 \tag{EC.1}
\]
The marked-node coordinate and source/Cartier rescaling are EN.1--3 of
\path{part3_e6_marked_node_cover_working.md}. At the selected cusp put
\(\xi=X-v\). The exact equation, including its full height, is
\[
 K=Y^2-\xi^2(\xi+w^2)
       =q\{q+v^2+D\xi\}.
 \tag{EC.2}
\]
Write \(a=v^2+D\xi\). On a sufficiently small common polydisc the
square root equal to \(a\) at \(K=0\) gives the literal inverse
\[
 q=\frac{2K}{a+\sqrt{a^2+4K}},\qquad
 T=v+q,\quad X=v+\xi.
 \tag{EC.3}
\]
Its denominator has value two at the origin. Thus \((\xi,Y,w)\)
are smooth coordinates on the actual source threefold; (EC.3) retains
the original elliptic-base function. In these coordinates the entire
selected rational fibre has normalization
\[
 f(S,w)=\bigl(S^2-w^2,S(S^2-w^2),w\bigr),\quad
 D_+=\{S=w\},\quad D_-=\{S=-w\},\quad
 Z=\frac{S-w}{S+w}.
 \tag{EC.4}
\]
The point at infinity is \(O\), and \(G\) has \(S=-1\). This is
the literal ordered germ used in F56, with \(w\) in place of its
ordered parameter. It is an immersion off \((S,w)=(0,0)\); its three
derivative minors generate \((S,w^2)\), and its only distinct double
pairs there are \((w,w),(-w,w)\). These statements follow directly
from (EC.4), with no base reparametrization.

The second selected critical value is
\[
 q_-=-4w^6/27.
 \tag{EC.5}
\]
Indeed the full discriminant is
\(-432T^2(T-v)(T-T_-)\), with \(v-T_-=4w^6/27\).
The unselected \(I_2\) remains at \(T=0\), outside a common small
disc about \(T=-1\). Hence the positive \(w\)-loop has exactly the
ordered two-node braid of F56, including its twelve positive half
twists and the actual boundary twist of its lifted once-punctured
torus. The orientation of that loop is preserved by the identity
\(w=\rho_{\rm F56}\), not inferred from the even original \(r(w)\).
In the original source parameter the exact identity is
\(r(w)-1/3=-w^4/(3D^2)\). Thus the present positive ordered loop
has local degree four over \(r-1/3\). Its actual source, line and
cap closing is the full periodic \(w\)-atlas below; no shorter
endpoint is assigned by this calculation.

For the full normal arrow put \(x=S-w,y=S+w\), so \(\xi=xy\), and
retain the nonzero covector
\[
 n_f=(3S^2-w^2)d\xi-2S\,dY+2w(S^2-w^2)dw.
\]
Restriction of the differential of (EC.2) to the normalization gives
\[
 dq=-\frac{xy}{v^2+Dxy}\,n_f.
 \tag{EC.6}
\]
In the two native fixed-edge charts \(q=Zc_0c_1\) and
\(q=Z^{-1}c'_0c'_1\), \(c'_i=Zc_i\), this is exactly
\[
 d(c_0c_1)=-\frac{y^2}{v^2+Dxy}n_f,
 \qquad d(c'_0c'_1)=-\frac{x^2}{v^2+Dxy}n_f.
 \tag{EC.7}
\]
The common factor has an analytic logarithm and square root on the
finite inner ball. These are the actual two Hom arrows of K42 SP.8,
with their full extending source unit. They also show explicitly why
the physical height cannot be replaced by its leading Tate term.

\par\medskip\noindent\textbf{EC.2. The actual Cartier lift and all selected period units.}\par\nopagebreak[4]

Put \(s_0=v(w^2+3)/4\) and
\(U=Y-T-DX/2,\ V_s=Y+T+DX/2\). Then
\(UV_s=X^2(X-s_0)\). With the actual Miller factors of EN.3 define
\[
 f_2=\frac{U}{X-s_0},\qquad
 f_3=\frac{s_0 XU}{s_0V_s-(Ds_0+2T)X},\qquad
 f_6=f_3^2f_{2,3G},
\]
\[
 B_6(p)=f_2(p)^6\frac{f_6(p-G)}{f_6(p)},\qquad
 B_{\rm sel}=\frac q{T^3}B_6.
 \tag{EC.8}
\]
These are coefficient arrows on the actual point line \(M_6\).
The point-line/tensor-line comparison is multiplication by
\(r_6=f_6^{-1}\); both its source and target factors occur in
(EC.8). At \(w=0,q=0,\xi=Y=0\), direct substitution gives
\[
 f_2=2,\quad f_3=3,\quad f_6=6,\quad
 B_6=64,\quad B_{\rm sel}/q=-64.
 \tag{EC.9}
\]
For the tangent factor at \(3G\), its cusp coordinates are
\((X,Y)=(-8/9,-1/27)\); its value at the cusp is \(2/3\),
which supplies the displayed value of \(f_6\). All horizontal
divisors of these factors avoid this cusp point. Thus these are
values of extending units, rather than cancellations of leading
terms at an indeterminate point.

There is a full selected-normalization check of its actual theta
frame. Put \(P=(1+w)/(1-w)\) and
\[
 S_6=w\frac{1+P^6}{1-P^6},\qquad
 c_6=w-S_6=\frac{(1+w)^6}{6+20w^2+6w^4}.
 \tag{EC.9a}
\]
The point \(6G\) is \(S=S_6\), with \(S_6(0)=-1/6\).
The exact divisor equations obey
\[
 S-S_6=c_6\,\frac{1-Z/P^6}{1-Z},\qquad
 f_6\big|_{q=0}=\frac{(S+1)^6}{S-S_6}.
 \tag{EC.9b}
\]
For the second identity both sides have the prescribed divisor
and the same leading coefficient \(S^5\) at \(O\); their ratio
is therefore the constant one on the normalization. Both branches
of the node have the same value
\[
 f_6\big|_{\rm node}=6+20w^2+6w^4,\qquad
 f_2\big|_{\rm node}=2,\qquad
 (B_{\rm sel}/q)\big|_{\rm node}=64/v^3.
 \tag{EC.9c}
\]
The reciprocal of the unit \(c_6\) is the coefficient arrow between
the displayed point-divisor and theta frames. It extends on the
entire parameter disc. This fixes the point-line theta gauge with
its full normalization constant, and shows explicitly that it
contributes no parameter winding.

Translation and inverse translation are also explicit here. For
\(+G\) use \(k=(Y-T)/X\), and for \(-G\) use \(k=(Y+T)/X\).
The respective target formulae are
\[
 X'=k^2-(3-2w^2)-X,\quad
 Y'=-T-kX'\quad\hbox{or}\quad Y'=T-kX'.
 \tag{EC.10}
\]
Since \(X=-1\) at the cusp, both are analytic there, and the two
maps are inverse by the actual chord law on their dense smooth
fibres. Their common fixed cusp therefore has the full line unit
(EC.9), including its constant sign. The scalar is
\(B_\lambda=\lambda B_{\rm sel}\), with the originally fixed
\(\lambda\); the factor \(-64\) is not absorbed into it.

On the selected normalization the marked source multiplier is
\[
 P_G=\frac{1+w}{1-w},\qquad g=P_G^{-1}.
 \tag{EC.11}
\]
The period and line charts give the exact commuting periods
\[
 H:(Z,V,q)\mapsto(\kappa qZ,\alpha V,q),\qquad
 B:(Z,V,q)\mapsto(gZ,\eta qV,q),\quad\eta=\lambda\beta.
 \tag{EC.12}
\]
Here \(\alpha\) and \(\beta\) are the coefficients of the actual
theta/Cartier trivialization of (EC.8), not new independently chosen
units. Its source point is \(6G\), so the marked theta-period
coefficient is \(P_G^6\); the actual rational function \(f_6\)
is elliptic and hence preserves that period on both copies. At the
node the invertible Hom coefficient is the one in (EC.9).

There is a joint domain for all the required logarithms. Substitute
\(\xi=w^2\widetilde X,Y=w^3\widetilde Y,q=w^6Q_c\) into (EC.2):
\[
 \widetilde Y^2=\widetilde X^2(\widetilde X+1)
       +Q_c(v^2+Dw^2\widetilde X)+w^6Q_c^2.
 \tag{EC.13}
\]
At \(w=0\) this is the marked compact pencil
\(\widetilde Y^2=\widetilde X^2(\widetilde X+1)+Q_c\).
Choose a sufficiently small fixed \(Q_c\)-disc containing just its
selected node. The inverse-j series, with its actual marked fourth
root, gives
\[
 \kappa=w^{-6}\widehat\kappa(Q_c,w),\qquad
 \widehat\kappa(0,w)=-v^2/64.
 \tag{EC.14}
\]
The factor \(\widehat\kappa\) and its inverse extend on this disc
times the full \(w\)-disc. The original section, viewed in the
weighted source chart at its section at infinity, has a formal
parameter
\[
 -\widetilde X(G)/\widetilde Y(G)
       =\frac{wv}{v+w^6Q_c},
 \tag{EC.14a}
\]
divisible by \(w\). Its denominator is a unit on the full
weighted bidisc. The marked Abel/theta point
therefore has an analytic formal-group logarithm tending to zero.
This proves the extension of \(P_G\), \(g\), and \(\alpha\) as
units in that same domain. Finally (EC.8)--(EC.9) supplies the
extending source and target Cartier units for \(\beta\); its
central value is \(-64\). Thus the based periodic logarithms of
all four units are literal, after retaining the sole angular factor
\(w^{-6}\) in \(\kappa\). Their winding matrix is
\(\operatorname{diag}(-6,0)\).

\par\medskip\noindent\textbf{EC.3. The unselected exterior and the actual finite cap.}\par\nopagebreak[4]

The original E6 source coordinate is
\[
 r=(3-2w^2)/D^2,\quad (x,y,t)=(D^{-2}X,D^{-3}Y,D^{-3}T).
 \tag{EC.15}
\]
On a small \(w\)-disc, \(D\) is a unit with a fixed analytic
logarithm, \(r\) lies in a disc about \(1/3\), and
\(s=1/4-r\) is a unit with value \(-1/12\). At the unselected
\(I_2\), the original coordinates satisfy
\[
 (y-t-x/2)(y+t+x/2)=x^2(x-s).
\]
The analytic map \(\widehat x=x\sqrt{x-s}\) identifies this with
\(UV=\widehat x^2\), retaining
\(t=(V-U-x)/2\). The two source resolution charts are
\(\widehat x=Uz,V=Uz^2\) and
\(\widehat x=V\zeta,U=V\zeta^2\), with \(z\zeta=1\).
Translation is the actual addition map of EU.1, including its
\(x\)-chart, the \(t\)-charts at \(\pm G\), and the zero-section
chart. The only denominator singled out in the exceptional
\(k=1/2\) chart is \(x=s\), which is a unit on this disc.
Consequently those maps and their inverses extend jointly here.

The divisor computation E68--69 gives the complete vertical divisor
\(\operatorname{div}_{\rm vert}B_6=3F_0-2F_\infty\).
The factor \(q/T^3\) in (EC.8) cancels it except at the selected
whole zero. At the parameter divisor \(w=0\), its value on any
ordinary reference section avoiding the stated horizontal divisors
is a unit. The resolved total source is smooth. The rational Hom
map and its inverse therefore extend at the remaining
codimension-two loci by their zero valuations at every height-one
prime. This carries the actual Hom coefficient on the entire
ordinary exterior; no parameter scalar is inserted.

At infinity put \(q_\infty=z^2\). The original quadratic root is
\[
 Y'^2=X'^3+z^2(rX'^2+X'+1),\quad G'=(0,z),
 \quad\mathfrak D(z,X',Y')=(-z,X',-Y').
 \tag{EC.16}
\]
With \(c=\sqrt{1+X'+rX'^2}\), the coordinates
\(U'=Y'-zc,V'=Y'+zc\) obey \(U'V'=X'^3\) and
\(z=(V'-U')/(2c)\). Use the three exact A2 charts of E70.
In the middle chart they are
\[
 X'=ab,\quad U'=a^2b,\quad V'=ab^2,\quad
 z=ab(b-a)/(2c),\quad Y'=ab(a+b)/2.
 \tag{EC.17}
\]
The full translation formula EU.5--6 is
\(a'=(b-a)F/H,b'=-aH\), with its actual analytic units
\(F,H\) and value-one denominators. These formulae, their inverse,
and the smooth group charts are uniform on this \(r\)-disc; none
requires \(r=0\). The half-division is the actual section
\(R=2G'-[2]^{-1}(3G')\); the inverse is the convergent formal-group
inverse at the identity. Hence \(2R=G'\) and \(\mathfrak D R=R\),
with its nonidentity central component retained.

Blow up the four literal deck-fixed points of E70. On that common
graph the two section equations at the zero-section blowup are
\(\tau'=zv_0\) and \(\tau'-\tau'(6G')=z(v_0-L_6(s,q_\infty))\),
where \(L_6(s,0)=-s\). The actual point-line arrow is
\[
 \frac{\tau'}{\tau'-\tau'(6G')}=
          \frac{v_0}{v_0-L_6(s,q_\infty)}.
 \tag{EC.18}
\]
It proves \(\beta^*\mathcal O(6G'-O)=\pi^*M_6\), including every
exceptional coefficient. On the root the full selected coefficient
is \(B_0=h_\infty z^{-2}B'_6\),
\(h_\infty=1-t_c(r)z^2\). This is invertible on every reduced-IV
component by E71's actual divisor computation.

For clarity retain the norm normalization. If a lift \(F\) of
\(t_{-R}\) has \(\mathfrak D F\mathfrak D^{-1}=c_FF\), then
\(c_F(-z)c_F(z)=1\), \(c_F(0)=1\). With its odd logarithm put
\(F_1=e^{(\log c_F)/2}F\). It commutes with \(\mathfrak D\).
The invariant unit \(\beta_0=B_0/F_1^2\) has its analytic square
root fixed at the full reference map. Then
\[
 F_2=\beta_0^{1/2}F_1,\quad
 A_0=(-1)^dF_2^n\mathfrak D,\quad A_0^2=B_0^n,
 \qquad B=R_\lambda^2 B_0,\quad A=R_\lambda^n A_0.
 \tag{EC.19}
\]
Here \(R_\lambda^2=\lambda\) is fixed, including its complex
phase. All source, line, character and norm arrows are equalities.

Under (EC.15), the root coordinates have the literal comparison
\[
 z_{\rm old}=D^{3/2}z_{\rm new},\quad
 X'_{\rm old}=D X'_{\rm new},\quad
 Y'_{\rm old}=D^{3/2}Y'_{\rm new}.
 \tag{EC.20}
\]
The square root uses the fixed logarithm of the unit \(D\).
The transformed root is exactly the equation in EN.3. The full
Miller weights are \(D^{-1},D^{-2},D^{-5}\), and
\(B_{6,\rm old}=D^{-6}B_{6,\rm new}\). Since
\((t-t_c)/t^3=D^6 q/T^3\), the complete suspension coefficient
is unchanged. The Cartier ratio (EC.18) carries both numerator
and denominator: explicitly
\(\tau'_{\rm old}=D^{-1/2}\tau'_{\rm new}\) and
\(\tau_{0,\rm old}=D^{-2}\tau_{0,\rm new}\), consistently with
\(\tau'=z\tau_0\). Also
\(1-t_{c,\rm old}z_{\rm old}^2=1-vz_{\rm new}^2\).
Thus (EC.20) closes literally over the \(w\)-disc;
no finite root rotation or scalar winding is concealed by the
source rescaling.

Here is the literal application of EU.11a--12 to this family. Work in
the current height \(T\), put \(q_\infty=T^{-1}\), and retain
\(t_c(w)=v=w^2-1\) and \(h_\infty=1-t_c(w)z^2\). On a fixed
smooth cap annulus choose the same two angular sectors and reference
ordinary-to-cap map as in that construction. Its actual
\(I_z\) is the source/root and Cartier map (EC.18) followed by
(EC.20); in particular \(I_{-z}=\mathfrak D I_z\). In the
rigidified extension-group Lie cover set
\[
 \begin{gathered}
 s_\alpha=\frac{\log_\alpha q_\infty}{2\pi i},\qquad
 z_\alpha=e^{\pi i s_\alpha},\qquad
 c_\lambda=\frac{n\log R_\lambda}{2\pi i}+\frac d2,\\
 f_\alpha=\operatorname{Log}_\alpha F_2
       =(-a_\alpha(R),\ell_\alpha),\qquad
 b_\alpha=\operatorname{Log}_\alpha B .
 \end{gathered}
 \tag{EC.21a}
\]
The source Abel integral \(a_\alpha(R)\) and the line logarithm
\(\ell_\alpha\) belong to the same extension group. The latter
retains the two \(r_6\) factors, \(h_\infty\), and the selected
square root of \(\beta_0\) in (EC.19). Thus exponentiation is
the full source/line arrow. It is not a choice of independent
scalar factors in unrelated Cartier charts.

Write \(s_\beta=s_\alpha+k_{\alpha\beta}\) and
\[
 f_\beta=\mathfrak D_*^{k_{\alpha\beta}}f_\alpha+
                         \Pi_{\alpha\beta}.
\]
Here \(\Pi_{\alpha\beta}\) is the actual transported source/line
period, with the integer coefficients fixed at the given reference
parameter \(w_*\). Retain also its reference suspension integers
\(j_{\alpha\beta}\) and full integral period vectors
\(\Lambda_{\alpha\beta}\). The coefficients are fixed while the
period functions themselves vary holomorphically with \(w\).
The ordinary-to-cap map is
\[
 \begin{split}
 \Phi_\alpha
   &=\operatorname{Exp}\bigl(
       n s_\alpha f_\alpha+s_\alpha c_\lambda\delta+
                         h_\alpha\bigr)I_{z_\alpha},\\
 h_\beta-\mathfrak D_*^{k_{\alpha\beta}}h_\alpha
   &=-n s_\alpha\Pi_{\alpha\beta}
          +j_{\alpha\beta}b_\beta+\Lambda_{\alpha\beta}.
 \end{split}
 \tag{EC.21b}
\]
This is EU.11c--d with \(r\) replaced by \(w\), with every displayed
root, Cartier, period and character map replaced by (EC.18)--(EC.20).
In particular its source displacement is
\(-n s_\alpha a_\alpha(R)\), and its scalar includes
\(R_\lambda^{n s_\alpha}\exp(\pi i d s_\alpha)\).

For completeness the correction in (EC.21b) is an actual
parameter-dependent map. Use \(I_z\) to put its Lie vectors into
the ordinary Lie bundle on the fixed annulus and its invariant
differential/rigidified-line frame. Extend the entire reference
\(h_\alpha(w_*)\) to \(h_\alpha^0\) in that frame. Subtract its
difference from the right side of (EC.21b), obtaining the additive
overlap vector \(E_{\alpha\beta}\), zero at \(w_*\).
For a fixed sector partition of unity put
\(k_\alpha=-\sum_\beta\chi_\beta E_{\alpha\beta}\).
Then \(k_\beta-k_\alpha=E_{\alpha\beta}\), and
\(\bar\partial k_\alpha\) is a single vector-valued form \(\omega\).
With a fixed cutoff on a slightly larger annulus use
\[
 v_{\rm CG}(q_\infty,w)=\frac1\pi\int
      \frac{\omega(\zeta,w)}{q_\infty-\zeta}\,dA(\zeta),
 \qquad h_\alpha=h_\alpha^0+k_\alpha-v_{\rm CG}.
 \tag{EC.21c}
\]
On the protected annulus \(\bar\partial v_{\rm CG}=\omega\).
This is holomorphic in \(w\), zero at \(w_*\), and solves the
specified overlap equation. Finite averaging is applied only to
these vector cochains; it preserves the fixed lattice cocycle,
character and complete reference map. Substitution in (EC.21b)
leaves, after subtracting the \(k_{\alpha\beta}\)-fold actual cap
action, exactly \(j_{\alpha\beta}b_\beta+\Lambda_{\alpha\beta}\).
Hence the full quotient arrows satisfy
\[
 \Phi_\beta=A^{k_{\alpha\beta}}B^{j_{\alpha\beta}}\Phi_\alpha,
 \qquad \Phi_\alpha B=B\Phi_\alpha .
 \tag{EC.21}
\]
Their inverse is \(I_z^{-1}\) followed by the negative Lie
argument. Both norm signs retain the fixed factor \(R_\lambda^n\)
and their actual finite character. All these maps are single valued
over the \(w\)-disc, so their continuation has identity closing.
This specifies the complete \(F_0\), rather than only its norm.
For the reflection used in EC.6, perform this exterior/root and
annular construction first on the disc \(\zeta=w^2\), with
reference \(\zeta_*=w_*^2\), and then pull it back by that exact
map. The source equation, \(G\), \(t_c=v\), \(D\), \(I_z\), and
all root/Cartier units above are holomorphic in \(\zeta\).
Choose the reference extensions \(h_\alpha^0\), period sections,
and Cauchy operator on that disc. The resulting complete
ordinary-to-cap arrow is literally even in \(w\) in these
original source coordinates. Its expression in the two ordered
selected Tate markings is then related by the exact reflection
arrow of EC.6; no additional \(w\)-dependent cochain is chosen.

The ordinary and cap quotients are proper uniformly on the closed
parameter disc. On their finite source charts away from the selected
zero the actual line maps and their inverses are bounded in fixed
Hermitian metrics. Choose the original fixed \(|R_\lambda|\) below
their common strict-contraction threshold. In the cap the literal
generator \(H=A^n\) satisfies \(H^2=B\), \(H^n=A\), and has scalar
factor \(R_\lambda\). Its norm shell is compact and has the
specified character; this proves freeness and properness for both
norm signs. The selected toric family on the compact punctured
w-circle has its finite union of protected period fundamental
regions. On the exterior of the selected disc the smooth charts,
the root resolution and these compact shells give a proper smooth
family over the entire w-disc, with the actual source/line boundary
marking (EC.21).

\par\medskip\noindent\textbf{EC.4. The real orbit product and the full fixed trace.}\par\nopagebreak[4]

On the small source region containing the selected pair the source
is the smooth threefold (EC.3). The section of its group defining
translation is in the smooth locus at the central cusp. Divide it
there by a sufficiently large integer, extend that division section
by the inverse function theorem, take its small-group logarithm,
and multiply the logarithm by that integer. Its exponential is the
actual section by uniqueness and the chord law. The invariant
vertical field, in the original source coordinates,
\[
 2Y\partial_X+\{3X^2+2(3-2w^2)X+DT\}\partial_Y
\]
extends across the cusp and vanishes there. Its indicated time maps
give the actual source translation isotopy. The nonzero line lift
\(B_{\rm sel}/q\) has the same extension-group construction; its
central scalar is the retained \(-64\). Its logarithm is chosen
at that value, not at one. This untwists the actual suspension on
the source product, with its outside collar fixed.

At the node use the exact increasing orbit coordinates of F9,
\(2b=r-|y_0|^2/r\), \(2b'=|x_0|^2r-r^{-1}\), \(db'/db>0\).
They give the real gamma lift on the common short arc; the branch
\(S=-w\) is lower and \(S=w\) upper in the marked coordinate
\(Z\). Their full-period closing is the identity on this ordered
\(w\)-cover. A nonzero gamma deck would change the norm of the
strictly contracting line and is absent from that identity.
All source-line logarithms on the outside and the cap extend on
the disc by EC.3. A horizontal field chosen there and on their
common boundary therefore yields the actual ambient orbit-body
product \(Y\times D_w\). Its fibre is the known orbit body of the
constructed E6 sphere. The F9 source-product arrow supplies this
extension also at the cusp, without extending the complex selected
toric quotient there.

Fractional source translations act on the normalization by the exact
formula
\[
 S\longmapsto
 \frac{(1+g_t)S+(1-g_t)w}
 {((1-g_t)/w)S+(1+g_t)},\qquad
 g_t=\exp(t\log g),
 \tag{EC.22}
\]
where \(g(0)=1\). If \(g_t=1+wh_t\), its difference from \(S\)
is \(h_t(S^2-w^2)/(2+h_t(w-S))\). The full derivative is
\(I+O(|S|+|w|)\) uniformly on a common sufficiently small ball.
Thus the written branch-fixing isotopy removes this actual source
shear. In the universal real gamma cover interpolate the actual
heights, within the common short arc, to
\(\operatorname{Re}(S\overline w)/\sqrt{|S|^2+|w|^2}\).
The two double points remain strictly ordered, so this is an
isotopy through embeddings. Its outer collar is unchanged.

Remove the star-shaped four-ball
\(f^{-1}(B^6_\epsilon)\) from \(\mathbf P^1\times D_w\).
The graph is the same actual embedded cobordism as F56 on the
normalization (EC.4). In the rounded ambient product ball
\(B^6_\epsilon\times[-2\eta\epsilon,2\eta\epsilon]\), its inner
link is literally \((f,\eta H)\), since \(|H|\le |w|\).
This is the F55 lifted-link map, with the same complex domain
orientation and branch-height order. Its outside maps have been
specified in EC.3 and (EC.22), rather than inferred from the cubic
germ alone.

\par\medskip\noindent\textbf{EC.5. The actual circle line, both caps and the reference pair.}\par\nopagebreak[4]

The pullback of the actual \(M_6\) to
\(\mathbf P^1\times D_w\) has degree zero on every normalization
fibre. Evaluation in its zero-section Cartier frame gives the
relative nonzero section and hence its literal product line.
Its value at the reference parameter fixes the remaining disc
unit; the full Hom isotopy of EC.4 is retained in that frame.
The normal determinant has the same covector and actual divisor
isomorphism as (EC.6):
\[
 L=\mathcal O(-D_+-D_-),\qquad
 E=\mathcal O(D_+)\oplus\mathcal O(-D_-),\qquad
 \det E\xrightarrow{\ (S-w)/(S+w)\ }\mathcal O.
 \tag{EC.23}
\]
In the positive spinor the physical nonzero-vector line is the
tautological line \(M\); its meridian-one Hermitian dual is \(U\).
The actual coefficient arrow of K42 is
\[
 (a,b)\longmapsto
 \frac1{\sqrt{|x|^2+|y|^2}}
 \begin{pmatrix}y&-\overline x\\x&\overline y\end{pmatrix}
 \overline{(a,b)}.
 \tag{EC.24}
\]
Its source-unit multiplier from (EC.7) has an extending logarithm
and square root. After matching the product of the two normal
coordinates, a residual gauge is \((c_0,c_1)\mapsto(gc_0,g^{-1}c_1)\).
Its unitary part gives the literal Hopf arrow
\(q_H\mapsto q_H g/|g|\), a physical-circle map. Its possible
parameter winding extends over \(D^3\times S^1\). This retains
all separate factors of the actual line/normal frame.

The outer cap map uses
\[
 A_w=\begin{pmatrix}1&-w\\1&w\end{pmatrix},\qquad
 (P,q_H,w)\mapsto
 \left(\frac{A_wPA_w^*}{\operatorname{tr}(A_wPA_w^*)},q_H,w\right).
 \tag{EC.25}
\]
Its inverse is the displayed formula with \(A_w^{-1}\). It holds
the actual Cox quaternion fixed, not its tangent derivative, and
is periodic on the full parameter circle. The determinant factor
\(2w\) is retained in (EC.24). The inner cap is the identical
projection-normal four-ball, with the extending unit in (EC.7).
The anti-linear dual map of K68 commutes with the SU(2) clutches
and extends over both normal discs; it carries \(M=U^{-1}\).

We next verify the global orbit hypotheses used by the Thom step.
The reference completed E6 member is the F41 sphere with datum
\((a,b,m,n)=(6,1,2,n)\), \(n=\pm1\), and the actual selected-zero
Hom section of EC.2. Its application to the present source is
literal: join the present marked value
\(\rho=w^2/(w^2-3)\) to \(\rho=-1\) by an arc whose interior
avoids \(0,1,-1,\tfrac12,-\tfrac12\). Choose a contractible
tubular neighbourhood and its ordered-node square root. The
Miller/Cartier formulas, (EC.18)--(EC.21), and their inverses give
the completed family on that tube. At the endpoint \(\rho=-1\),
\(r=0\), the K94 unselected-III extension is the actual E6-B
source of F41. The fixed sufficiently small scalar makes this a
proper compact family along the arc, carrying the reference
cochain and character. Its Ehresmann transport identifies the
present reference member \(X\) with that sphere. No loop or
additional parameter return is inserted by this interval map.

The physical circle is free on the ordinary part. In the quadratic
cap the actual generator \(H=A^n\) satisfies \(H^2=B\) and
\(H^n=A\), and its strictly contracting norm prevents a nonzero
\(H\)-power from agreeing with a unit physical-circle rotation.
The standard cusp fan has edge directions
\(\pm e_1,\pm e_2,\pm(e_1-e_2)\). The determinant with the
physical cocharacter \(e_2\) is \(\pm1\) off the vertical
edges, whose period quotient is one fixed curve \(F=\mathbf P^1\).
The two charts
\[
 c_0'=Zc_0,\quad c_1'=Zc_1,\quad Z'=Z^{-1},
 \qquad q=Zc_0c_1=Z'c_0'c_1'
 \tag{EC.25a}
\]
have weights \((-1,+1,0)\). Thus the fixed normal bundle and
semifree action are precisely the displayed local input of
F7--8, including the physical circle generator.

Put \(E=X\setminus\operatorname{int}(F\times D^4)\) and
\(P=E/S^1\). Since \(X\cong S^6\), Alexander duality gives the
integral homology of \(S^3\) for \(E\), and codimension four
general position gives \(\pi_1(E)=0\). Conjugating the negative
normal line in (EC.25a) gives
\(\mathcal O(1)\oplus\mathcal O(-1)\). Its clutch
\(\operatorname{diag}(Z,Z^{-1})\) bounds the SU(2) disc with
coordinates \((r e^{i\theta},\sqrt{1-r^2})\); adjoining a path
from its centre to the identity gives the actual unitary
trivialization. In that trivialization the boundary quotient is
the product Hopf map
\[
 \partial E=S^2_F\times S^3
       \longrightarrow S^2_F\times S^2_{\rm mer}.
 \tag{EC.25b}
\]
Mayer--Vietoris for \(X=E\cup(F\times D^4)\) identifies
\(H_3(\partial E)\to H_3(E)\) as an isomorphism. Its restriction
to a normal \(S^3\) is the Hopf bundle, whose fibre integration
is the primitive meridian Euler generator. Naturality in the
integral Gysin sequence consequently makes
\(H^3(E)\to H^2(P)=\mathbb Z e\) an isomorphism. The same
sequence gives \(H^3(P)=H^4(P)=0\); the homotopy sequence gives
\(\pi_1(P)=0\). The Euler class restricts to the meridian
generator and has zero \(S^2_F\) coefficient, by (EC.25b).

For \(Y=P\cup(F\times D^3)\), the integral Mayer--Vietoris
map in degree two, with boundary basis \((F,\mathrm{meridian})\)
and target \((P,\mathrm{tube})\), is
\[
 \begin{pmatrix}0&1\\-1&0\end{pmatrix}.
 \tag{EC.25c}
\]
Its determinant is one. The boundary \(H_4=\mathbb Z\) maps to
zero and supplies \(H_5(Y)=\mathbb Z\); the other intermediate
groups vanish. Van Kampen gives \(\pi_1(Y)=0\).
The verified \(\Theta_5=0\) and Haefliger applications in F7--8
therefore identify the orbit pair as \((S^5,\text{unknotted }S^2)\).
On its free complement the actual Hopf bundle has Euler
generator. Its comparison with the fixed normal product differs
on \(S^2\times S^2\) by a circle gauge; that gauge has a global
logarithm because \(H^1(S^2\times S^2;\mathbb Z)=0\), and the
logarithm extends over the collar. These are the actual
equivariant reference maps used in (EC.24)--(EC.25).

The ambient meridian line extends on the excision complement by
the same integral Thom map as F56: the product orbit body has
\(H^2=H^3=0\), and the properly embedded connected normalization
has meridian evaluation \(H^2\cong\mathbb Z\). The actual outer
line has the specified value, and its ratio gauge extends over
the written outer cap. Thus the extra normal-sphere boundary is
the actual bundle over
\((S^2\times D^2)\cup(D^3\times S^1)=S^4\), with both caps and
their normal clutch fixed by (EC.24)--(EC.25).

Consequently every step of the F56/K42 excision uses the same
oriented link, coefficient matrix and cap morphisms. The relative
second Chern number is \(-1\), computed from (EC.24) with cap
boundary orientation negative to \(\mathbf P^1\times S^1\).
The normal-sphere pair is \((1,4)\) for \(U\); the inner lifted
link pair is \((5,-4)\) in its retained convention. Their actual
boundary signs and the Jin parameter reversal give the reference
\(U\)-pair \((6,0)\). The Hermitian dual map gives the physical
pair \((-6,0)\). The K68 first-to-second piece clutch is
\(\lambda_M(F)=-\Sigma_{F^{-1}}=\Sigma_F\); the return is the
one obtained from the delta-averaged product connection, identical
on the selected tube, so this is the actual clutched sphere. Its
reference class and native coordinates are therefore the K71 seed
\(-1\) and \((0,-6,-1)\), respectively. These conclusions use the
complete exterior and cap extension above.

\par\medskip\noindent\textbf{EC.6. The actual vertical seeds and complete fanwise trace comparison.}\par\nopagebreak[4]

First retain the complete vertical seed; the initial seed alone does
not specify a different vertical star. On \(T_j\) the literal
characters and periods give
\[
 V_j=Z^{-j}V,\quad H_j=HB^j,\quad B_j=B,\quad
 Q_j=\kappa qg^j,
\]
\[
 \alpha_j=\alpha\eta^j\kappa^{-j}g^{-j^2},
 \qquad\eta_j=\eta g^{-j}.
 \tag{EC.26a}
\]
For each fixed integer \(j\), the exact source interpolation is
\(Q_t=Q\exp(tj\log g)\). Near its node \(xy=Q_t\), the vector
field \(\partial_t+j\log g\,y\partial_y\) is tangent and keeps
\(q,w\) fixed. On the outside torus use the marked-period field
\(u+\tau_t b\mapsto u+\tau_0b\). Their difference on the
overlap is vertical. A partition of unity in this affine space
of horizontal fields, fixed on the two protected collars, supplies
the actual periodic flow. It transports the embedded normalization
graph by the ordered real-height isotopy of EC.4. The spinor arrow
is multiplication by \(Z^j\):
\[
 E\otimes\mathcal O(j(D_+-D_-))\longrightarrow E.
\]
Its two Cartier generators \(x^{-j-1}y^j,x^{-j}y^{j+1}\) go
exactly to \(1/x,y\). Its adjoint and the product cap (EC.25)
therefore carry the full normal frame, including (EC.7).

In the full real period coefficients \((u,b,d,\gamma)\), the
actual outer arrow is
\(u_j=u,b_j=b,\gamma_j=\gamma-jb,d_j=d-ju\).
The q-deck is \((u,d)\mapsto(u-b,d-\gamma)\), and the positive
w-deck is \((u,d)\mapsto(u+6b,d)\). These follow from (EC.12)
with every actual unit logarithm retained. In those coefficient
frames the connections are
\[
 A_0=dd+\alpha_q d\gamma,\qquad
 A_j=dd_j+\alpha_q d\gamma_j+6j\theta\,db,
\]
\[
 A_j-A_0=-j\{du+(\alpha_q-6\theta)db\}.
 \tag{EC.26b}
\]
The differentials of the zero-winding logarithms occur in the
definition of \(d\); using these coefficients retains them.
The only new integral theta-clutch is the dual-line factor
\(\mathcal T^{6j}\), where the coefficient of \(\mathcal T\)
changes by \(e^{2\pi i b}\). Its curvature is
\(\Xi=d\theta\wedge db\), with \(\Xi^2=0\).

Here the exact relative extension test of K66 applies to the
present source, with the following concrete maps. The graph map
\(b\) is constant on the attaching vanishing circle and extends
over its two-handle; it is made constant near the fixed graph by
a supported real lift and then extends over the product cap.
The integral boundary has
\(H_1=\mathbb Z\langle\alpha_q,\theta,b,\gamma\rangle\);
the q commutator kills \(u\) primitively. The restrictions from
the completed piece contain precisely \(\theta,b,\gamma\).
The remaining q-angle obstruction vanishes on the literal disc
\[
 (q/Z_*,Z_*,1)\mapsto(qZ_*^{j-1},Z_*,Z_*^{-j}),
 \qquad |Z_*|=1,\quad Z_*\text{ near }-1,\quad\gamma=0.
 \tag{EC.26c}
\]
For small \(w\), the marked Cartier sections have Tate points
near one, so this disc avoids them and the fixed tube. Its actual
line arrow extends with no zero. Subtract the integral one-cochain
from the logarithms of the ratio line, solve its real Cech
one-cocycle by a partition of unity, and correct its boundary value
by the extending \(\theta,b,\gamma\) characters and the remaining
real logarithm. This gives the prescribed relative isomorphism
\(U_j\to U_0\otimes\mathcal T^{6j}\), not just its Chern class.

The full tangent comparison is the source flow just written and
its actual positive triangular graph shear. The source-frame
matrix is upper triangular after ordering its coframe as in
K63 SQ.4; its strictly triangular curvature has zero trace-square.
The base connection and its interpolation depend on only the
two boundary angles and one collar variable, so its curvature
square vanishes as well. The relative \(p_1\)-pairing against
\(\Xi\) is zero: a regular b-level outside the surgery is the
actual source three-level times the gamma circle, with the source
four-form restricting to zero and the just specified boundary
transgression zero. This retains the full E6 source connection.

The relative Euler cube has an affine interior term because
\(\Xi^2=0\). On the outward collar the exact difference (EC.26b)
has the same fifth-form coefficient \(18j^2\) as K66 EC.29;
the actual piece-to-exterior interpolation runs from one to zero,
so its physical contribution is \(-18j^2\). Thus the physical
cube has the form \(-6+Lj-18j^2\), and its Pontryagin pairing
stays zero. The two literal endpoint evaluations fix \(L\).
The E6 equation, the section, and all Miller factors depend only
on \(w^2\). The full reflection line arrow can be read directly
from the retained theta frame. One has
\[
 S_6(-w)=S_6(w),\qquad c_6(-w)=P^{-6}c_6(w),\qquad
 f_\theta(Z^{-1},P^{-6})=P^6f_\theta(Z,P^6).
\]
The last identity holds for the full theta product at the same
period \(Q\): use
\(\Theta(Z^{-1})=-Z^{-1}\Theta(Z)\) in its numerator and
denominator. In the first two identities \(P=P_G(0,w)\).
In the full theta identity and the following period arrows,
\(P=P_G(q,w)\) is the actual extending marked lift. The source
identity reverses the marked invariant differential and fixes
\(O\), so it sends this lift to \(P^{-1}\); the integral
Tate reindex is zero on the formal neighbourhood of \(O\)
specified by (EC.14a). Therefore the source/Cartier identity first has
the literal ordered-coordinate arrow
\[
 I:(Z,V)\mapsto(Z^{-1},P^6V).
\]
The actual target source period has coefficient
\(\alpha_{\rm target}=P^{-6}\). Composing that period with \(I\)
gives exactly
\[
 H_{\rm target} I:(Z,V)\mapsto(Q/Z,V).
\]
This retains the full height-one source-period reindex and all
theta factors. This composed half-turn arrow commutes with the
exterior and the finite cap (EC.19)--(EC.21), whose coefficients
are pulled back from the same source disc. This arrow sends
\(T_0\) to its literal reflected \(T_1\). It carries (EC.24)--(EC.25)
and the physical line, so both reference cubes are \(-6\), with
Pontryagin pairing zero. Hence \(L=18\), and on every \(T_j\)
the actual physical pair, class and native pair are
\[
 (k_M,\ell_M)=(-6-18j(j-1),0),\quad
 A_M=-1-3j(j-1),\quad(m_{\rm native},n_{\rm native})
       =(0,-6-18j(j-1)).
 \tag{EC.26d}
\]
All scalar/root and higher-unit terms in these maps are the ones
of EC.2--3. This establishes the complete K71 seed by the actual
E6 line and collar maps.

For a chosen lattice-basis fan \(P\), put \(R=P^{-1}\) and use
the literal characters \(Z'=Z^{R_{11}}V^{R_{12}}\),
\(V'=Z^{R_{21}}V^{R_{22}}\), with \(q\) unchanged. The new period
basis is \(H'_j=H^{P_{1j}}B^{P_{2j}}\). Its full coefficients are
\[
 A_{ij}=(\kappa^{P_{1j}}g^{P_{2j}})^{R_{i1}}
              (\alpha^{P_{1j}}\eta^{P_{2j}})^{R_{i2}},
 \qquad B_P=-6P^{-1}e_1e_1^tP.
 \tag{EC.26}
\]
The q-exponents are \(RP=I\). Every cone map and its inverse is
the corresponding integral semigroup map, carrying all full unit
powers. The outside arrow is its inverse character map in the
original source and Cartier coordinates, so the cap squares of
EC.3 are unchanged.

The three exact edge transitions, in the two vertex charts, are
\[
 (u,v_2,w_2)=
 (A_{11}A_{21}bc,A_{21}^{-1}ab,A_{11}^{-1}b^{-1}),
\]
\[
 (A_{12}A_{22}bc,A_{22}^{-1}c^{-1},A_{12}^{-1}ac),
 \qquad(a^{-1},ab,ac).
 \tag{EC.27}
\]
Their full normal matrices are
\(\operatorname{diag}(A_{22}/A_{21},A_{11}A_{21})\),
\(\operatorname{diag}(A_{11}/A_{12},A_{12}A_{22})\), and
\(\operatorname{diag}(A_{22},A_{11})\). Their determinant winding
is \(-6\). Retain their first winding numbers
\(l=(B_{22}-B_{21},B_{11}-B_{12},B_{22})\).
The complete unit contractions are the GF.9 maps
\(H_d=(\sin\vartheta e^{id\theta}+j\cos\vartheta)
(\sin\vartheta+j\cos\vartheta)^{-1}\), with left \(H_{-3}\)
and right \(H_{l_i+3}\), multiplied by the based logarithms
specified in EC.2 and their actual radial cutoffs. The exact tube
identity \(zD_{i,*}K_i=zD_i\) matches both vertex collars. Normal
rescaling and the cutoff flow of its velocity carry the full
higher normal jets and preserve their overlap maps.

Thus the actual ordinary/native boundary comparison is precisely
GF.12, with left degree \(-6\), right degree zero, and inverse
cylinder correction \(L^6\). Physical weights, odd Spin endpoints,
the capped edge traces and all nine prism equations are the same
integer equations of K71, derived from these actual matrices.
The local Cox ruling change is the same map on its rank-one matrix,
with common link \((h,q_H)\mapsto(q_H^{-1}hq_H,q_H)\); it is
independent of the source parameter and intertwines every full
coefficient and closing map in (EC.26). Its final native reset
retains the physical circle and central generator orientation.
Hence finite Euclidean words construct the same complete K71
algorithm on the present E6 family, starting from the actual seed
of EC.5. This proves equality with the K71 actual return on each
specified fan by these source/exterior/cap/Cox arrows. It assigns
no return to any other E6 puncture.

\par\medskip\noindent\textbf{EC.7. Audit scope and exact certificate.}\par\nopagebreak[4]

The certificate \path{code/part3_fourfold/e6_selected_cusp.py} checks
the full source-height equation and inverse polynomial relation,
normalization and normal covector, critical separation, actual
Miller values, source translation at the cusp, original/root
coordinate scaling, the full selected scalar cancellation, and
the exact period-conjugation and representative Cox identities.
The exterior, integral cochain, both normal-cap and physical-clutch
arguments are written above for geometric audit. The source
equation alone is not used to replace those arrows.
\end{proof}
